\paragraph{共动缺陷证书的横向可见窗与有限覆盖禁入}\label{par:xi-comoving-transverse-window-band-exclusion}
沿用 Cayley 紧化 \(\mathcal{C}:\mathbb{H}\to\mathbb{D}\)（\eqref{eq:app-cayley-disk}）与共动盘内像
\(
w_{\rho,x_0}=\mathcal{C}\bigl((\gamma-x_0)+\mathrm{i}\delta\bigr)
\)
（\eqref{eq:xi-comoving-disk-image}）。
记可见窗半长度 \(X(\varrho,\delta)\) 为 \eqref{eq:xi-single-defect-X} 所定义的函数；由 Cayley 模长闭式 \eqref{eq:xi-cayley-modulus-delta} 可知：对任意 \(0<\varrho<1\)、\(\delta\in(0,1)\) 与 \(x\in\RR\)，成立等价关系
\begin{equation}\label{eq:xi-transverse-visible-window-equivalence}
\boxed{\;
\bigl|\mathcal{C}(x+\mathrm{i}\delta)\bigr|<\varrho
\quad\Longleftrightarrow\quad
|x|<X(\varrho,\delta)\; }.
\end{equation}

\begin{theorem}[共动缺陷证书诱导的横向可见窗：局部离线零点禁入]\label{thm:xi-comoving-defect-transverse-visible-window}
给定 \(0<\varrho<1\)、\(\varepsilon\ge 0\) 并定义排斥半径
\[
r_\ast:=\varrho e^{-\varepsilon}\in(0,\varrho].
\]
取任意 \(\delta_0\in(0,1)\) 与 \(x_0\in\RR\)。
若满足
\[
D_{\zeta,x_0}(\varrho)\le \varepsilon,
\]
则对任意离临界线零点 \(\rho=\beta+\mathrm{i}\gamma\)（\(\beta<\tfrac12\)，并记 \(\delta:=\tfrac12-\beta>0\)）只要 \(\delta\ge\delta_0\)，必有
\begin{equation}\label{eq:xi-transverse-window-separation}
\boxed{\;
|\gamma-x_0|\ \ge\ X(r_\ast,\delta_0)\; }.
\end{equation}
因此当 \(X(r_\ast,\delta_0)>0\)（等价于 \(r_\ast>\frac{1-\delta_0}{1+\delta_0}\)）时，不存在同时满足
\[
\beta\le\tfrac12-\delta_0,
\qquad
|\gamma-x_0|<X(r_\ast,\delta_0)
\]
的离线零点。
\end{theorem}

\begin{proof}
反设存在离线零点 \(\rho=\beta+\mathrm{i}\gamma\) 满足 \(\delta=\tfrac12-\beta\ge\delta_0\) 且
\(
|\gamma-x_0|<X(r_\ast,\delta_0)
\)。
由 \eqref{eq:xi-transverse-visible-window-equivalence}（取 \(\varrho=r_\ast,\ \delta=\delta_0\)）得到
\[
\bigl|\mathcal{C}((\gamma-x_0)+\mathrm{i}\delta_0)\bigr|<r_\ast.
\]
另一方面，\(\Xi_{\zeta}\) 的零点位于临界带 \(0<\Re(\rho)<1\)，故 \(0<\delta<1\)。
由 \eqref{eq:xi-cayley-modulus-delta} 的闭式分式表示可见：对固定 \(x\in\RR\)，映射
\(
\delta\mapsto |\mathcal{C}(x+\mathrm{i}\delta)|^2
=\frac{x^2+(1-\delta)^2}{x^2+(1+\delta)^2}
\)
在 \(\delta\in(0,1)\) 上严格递减；
因此由 \(\delta\ge\delta_0\) 得
\[
\bigl|w_{\rho,x_0}\bigr|
=\bigl|\mathcal{C}((\gamma-x_0)+\mathrm{i}\delta)\bigr|
\le \bigl|\mathcal{C}((\gamma-x_0)+\mathrm{i}\delta_0)\bigr|
<r_\ast\le\varrho.
\]
此时 \(w_{\rho,x_0}\) 为 \(F_{\zeta,x_0}\) 的盘内零点，命题 \ref{prop:app-jensen-single-zero-lower-bound}（单零点 Jensen 下界）给出
\[
D_{\zeta,x_0}(\varrho)\ \ge\ \log\!\Bigl(\frac{\varrho}{|w_{\rho,x_0}|}\Bigr)\ >\ \log\!\Bigl(\frac{\varrho}{r_\ast}\Bigr)=\varepsilon,
\]
与 \(D_{\zeta,x_0}(\varrho)\le\varepsilon\) 矛盾。
故 \eqref{eq:xi-transverse-window-separation} 成立。证毕。
\end{proof}

\begin{corollary}[有限共动中心的覆盖式带状禁入：有限核验替代连续上确界]\label{cor:xi-comoving-finite-centers-band-exclusion}
固定 \(\delta_0\in(0,1)\)、\(0<\varrho<1\)、\(\varepsilon\ge 0\)，并令 \(r_\ast=\varrho e^{-\varepsilon}\)、\(H:=X(r_\ast,\delta_0)\)。
假设 \(H>0\)。
给定高度窗 \(T>0\)，取等距中心集
\[
\mathcal X:=\{-T+H+2Hj:\ j=0,1,\dots,N\},
\qquad
N:=\left\lceil\frac{T-H}{H}\right\rceil.
\]
若对所有 \(x_0\in\mathcal X\) 均有
\[
D_{\zeta,x_0}(\varrho)\le \varepsilon,
\]
则不存在离线零点 \(\rho=\beta+\mathrm{i}\gamma\) 同时满足
\[
\beta\le\tfrac12-\delta_0,
\qquad
|\gamma|\le T-H.
\]
\end{corollary}

\begin{proof}
由构造可知相邻中心间距为 \(2H\)，并且并集
\(
\bigcup_{x_0\in\mathcal X}(x_0-H,\ x_0+H)
\)
覆盖区间 \((-T,\ T)\)；
因此对任意 \(\gamma\) 满足 \(|\gamma|\le T-H\)，存在 \(x_0\in\mathcal X\) 使 \(|\gamma-x_0|<H\)。
若存在离线零点且 \(\delta\ge\delta_0\)，则与定理 \ref{thm:xi-comoving-defect-transverse-visible-window} 矛盾。证毕。
\end{proof}

\begin{theorem}[黄金分辨率链上的可见窗增长与覆盖复杂度的标度律]\label{thm:xi-golden-chain-visible-window-growth}
令 \(\varphi=\frac{1+\sqrt5}{2}\)，并取黄金分辨率链
\[
\varrho_m:=\tanh\!\Bigl(\frac12 m\log\varphi\Bigr)\qquad(m\to\infty)
\]
（\eqref{eq:xi-golden-resolution-chain-rho}）。
令 \(\varepsilon_m\ge 0\) 满足 \(\varepsilon_m=o(\varphi^{-m})\)，并置
\(
r_{\ast,m}:=\varrho_m e^{-\varepsilon_m}
\)。
对任意固定 \(\delta_0\in(0,1)\) 定义
\[
H_m:=X(r_{\ast,m},\delta_0).
\]
则当 \(m\to\infty\) 时成立渐近律
\begin{equation}\label{eq:xi-golden-chain-visible-window-growth}
\boxed{\;
H_m=\sqrt{\delta_0}\,\varphi^{m/2}\,(1+o(1))\; }.
\end{equation}
并且在高度窗 \([{-}T,T]\) 内排除所有满足 \(\beta\le\tfrac12-\delta_0\) 的离线零点，只需在 \(O(T/H_m)\) 个共动中心上核验不等式 \(D_{\zeta,x_0}(\varrho_m)\le\varepsilon_m\)。
因此若取 \(m=2\lceil\log_\varphi T\rceil+O(1)\)，则所需核验中心数可取为 \(O(\delta_0^{-1/2})\)，并在 \(T\to\infty\) 时保持有界于该量级。
\end{theorem}

\begin{proof}
由 \eqref{eq:xi-golden-resolution-chain-rho-closed} 得
\(
1-\varrho_m^2=4\varphi^{-m}+O(\varphi^{-2m})
\)。
又因 \(\varepsilon_m=o(\varphi^{-m})\)，有
\(
r_{\ast,m}^2=\varrho_m^2e^{-2\varepsilon_m}=\varrho_m^2+o(\varphi^{-m})
\)，从而
\[
1-r_{\ast,m}^2
=(1-\varrho_m^2)+o(\varphi^{-m})
=4\varphi^{-m}\,(1+o(1)).
\]
将其代入 \eqref{eq:xi-single-defect-X}（取 \(\varrho=r_{\ast,m},\ \delta=\delta_0\)）并注意当 \(m\) 充分大时分子趋于 \(4\delta_0>0\)，得到
\[
H_m^2
=\frac{r_{\ast,m}^2(1+\delta_0)^2-(1-\delta_0)^2}{1-r_{\ast,m}^2}
=\frac{4\delta_0+o(1)}{4\varphi^{-m}(1+o(1))}
=\delta_0\,\varphi^{m}\,(1+o(1)),
\]
从而得到 \eqref{eq:xi-golden-chain-visible-window-growth}。
覆盖中心数的结论由推论 \ref{cor:xi-comoving-finite-centers-band-exclusion} 的显式构造给出：用步长 \(2H_m\) 的等距中心覆盖 \((-T,T)\) 即可，中心数为 \(O(T/H_m)\)。
最后取 \(m=2\lceil\log_\varphi T\rceil+O(1)\) 使 \(\varphi^{m/2}\asymp T\)，代入即得中心数为 \(O(\delta_0^{-1/2})\)。证毕。
\end{proof}

\begin{proposition}[有限共动采样的不可避免盲区：纯几何覆盖下界]\label{prop:xi-comoving-finite-centers-blindspot}
固定 \(\delta_0\in(0,1)\) 与 \(0<\varrho<1\)，并令
\[
H_0:=X(\varrho,\delta_0).
\]
若 \(H_0>0\)，则对任意有限中心集 \(\{x_1,\dots,x_K\}\subset\RR\)，存在 \(\gamma\in\RR\) 使
\[
|\gamma-x_j|\ge H_0\qquad(1\le j\le K),
\]
从而点 \(\rho=\bigl(\tfrac12-\delta_0\bigr)+\mathrm{i}\gamma\) 的共动盘内像满足
\[
\bigl|w_{\rho,x_j}\bigr|\ge \varrho\qquad(1\le j\le K),
\]
并且该零点在半径 \(\varrho\) 处对所有采样中心均不可见（不进入 Jensen 求和式的 \(|w|<\varrho\) 求和域）。
因此，任何仅依赖有限数据族 \(\{D_{\zeta,x_j}(\varrho)\}_{j=1}^{K}\) 的核验机制无法在逻辑上排除该类零点的存在。
\par
特别地，若目标是在 \(|\gamma|\le T\) 内排除所有满足 \(\beta\le\tfrac12-\delta_0\) 的离线零点，则仅凭固定半径 \(\varrho\) 的有限中心采样，必要条件为
\[
K\ \ge\ \left\lceil\frac{T}{H_0}\right\rceil.
\]
\end{proposition}

\begin{proof}
由 \(H_0>0\) 可知区间 \((x_j-H_0,\ x_j+H_0)\) 的并集为有限个有界开区间之并，故其补集非空；取
\(
\gamma\in\RR\setminus\bigcup_{j=1}^{K}(x_j-H_0,\ x_j+H_0)
\)
即得 \(|\gamma-x_j|\ge H_0\)（\(\forall j\)）。
由 \eqref{eq:xi-transverse-visible-window-equivalence}（取 \(\varrho=\varrho,\ \delta=\delta_0\)）的对偶形式可得
\(
|\gamma-x_j|\ge H_0\Rightarrow |\mathcal{C}((\gamma-x_j)+\mathrm{i}\delta_0)|\ge\varrho
\)，故 \(|w_{\rho,x_j}|\ge\varrho\)。
于是该零点不会对 \(D_{\zeta,x_j}(\varrho)\) 的 Jensen 求和项产生贡献，因而有限采样数据无法导出矛盾。
最后一条为区间覆盖下界：要使对所有 \(|\gamma|\le T\) 存在某个中心满足 \(|\gamma-x_j|<H_0\)，则 \(K\) 个半径 \(H_0\) 的区间必须覆盖长度 \(2T\) 的区间 \([-T,T]\)，从而 \(2KH_0\ge 2T\) 并得 \(K\ge\lceil T/H_0\rceil\)。证毕。
\end{proof}

\paragraph{Poisson--Cauchy 粗粒化下相对熵的四阶主项与二阶矩锁定}\label{par:xi-poisson-cauchy-kl-fourth-order}
令 Poisson/Cauchy 核
\[
P_t(x):=\frac{1}{\pi}\frac{t}{t^2+x^2}\qquad(t>0),
\]
并对概率密度 \(h,g\) 定义粗粒化 \(h_t:=P_t*h\)、\(g_t:=P_t*g\)。
下述结论给出 \(t\to\infty\) 时相对熵的精确主项，并表明四阶主项由二阶矩差唯一锁定。

\begin{theorem}[Poisson--Cauchy 粗粒化下相对熵的四阶渐近：主项系数的普适性]\label{thm:xi-poisson-cauchy-kl-fourth-order-universality}
设 \(h,g\) 为 \(\RR\) 上概率密度，满足：
\begin{enumerate}
  \item 同均值规范化 \(\int_{\RR}x\,h(x)\dd x=\int_{\RR}x\,g(x)\dd x=0\)；
  \item 三阶绝对矩有限 \(\int_{\RR}|x|^3h(x)\dd x+\int_{\RR}|x|^3g(x)\dd x<\infty\)；
  \item \(h,g\) 在紧集上有界（用于控制一致余项）。
\end{enumerate}
记二阶矩差
\[
\Delta_2:=\int_{\RR}x^2\bigl(h(x)-g(x)\bigr)\dd x.
\]
则当 \(t\to\infty\) 时成立精确渐近
\begin{equation}\label{eq:xi-poisson-cauchy-kl-fourth-order}
\boxed{\;
D_{\mathrm{KL}}(h_t\mid g_t)
=\frac{\Delta_2^2}{8}\,t^{-4}+O(t^{-5})\; }.
\end{equation}
特别地，主项系数 \(1/8\) 与 \(h,g\) 的具体形状无关，为 Poisson--Cauchy 粗粒化在一维上的普适常数。
\end{theorem}

\begin{proof}[证明要点]
由定义 \(h_t(x)=\int_{\RR}P_t(x-y)h(y)\dd y\)。
对固定 \(x\) 以 \(y\) 为小参数（尺度 \(t\)）对核 \(P_t(x-y)\) 作 Taylor 展开，并利用 \(\int h=1\)、\(\int yh(y)\dd y=0\) 得到
\[
h_t(x)=P_t(x)+\frac{m_2(h)}{2}\,P_t''(x)+O(t^{-4}),
\qquad
g_t(x)=P_t(x)+\frac{m_2(g)}{2}\,P_t''(x)+O(t^{-4}),
\]
其中 \(m_2(h)=\int y^2h(y)\dd y\)，余项由三阶绝对矩与 \(\|P_t^{(3)}\|_{L^1}\asymp t^{-3}\) 控制。
因此
\[
\eta_t(x):=h_t(x)-g_t(x)=\frac{\Delta_2}{2}\,P_t''(x)+O(t^{-4}).
\]
又由于 \(g_t=P_t+O(t^{-3})\) 且 \(\eta_t/g_t=O(t^{-2})\) 一致成立，对 \((1+u)\log(1+u)\) 作二阶展开得到
\[
D_{\mathrm{KL}}(h_t\mid g_t)
=\frac12\int_{\RR}\frac{\eta_t(x)^2}{g_t(x)}\dd x+O(t^{-5})
=\frac{\Delta_2^2}{8}\int_{\RR}\frac{(P_t''(x))^2}{P_t(x)}\dd x+O(t^{-5}).
\]
最后由 Poisson--Cauchy 能量恒等式 \eqref{eq:cdim-poisson-pt-second-energy} 得
\(
\int_{\RR}\frac{(P_t'')^2}{P_t}\dd x=t^{-4}
\)，代回即得 \eqref{eq:xi-poisson-cauchy-kl-fourth-order}。证毕。
\end{proof}

\begin{corollary}[\texorpdfstring{$o(t^{-4})$}{o(t^{-4})} 迫使二阶矩锁定]\label{cor:xi-poisson-cauchy-kl-moment-locking}
在定理 \ref{thm:xi-poisson-cauchy-kl-fourth-order-universality} 的条件下，若
\[
D_{\mathrm{KL}}(h_t\mid g_t)=o(t^{-4})\qquad(t\to\infty),
\]
则必有 \(\Delta_2=0\)，亦即 \(\int x^2h=\int x^2g\)。
\end{corollary}

\begin{proof}
由 \eqref{eq:xi-poisson-cauchy-kl-fourth-order} 取极限得
\(
t^4D_{\mathrm{KL}}(h_t\mid g_t)\to \Delta_2^2/8
\)。
若左端趋于 \(0\)，则 \(\Delta_2=0\)。证毕。
\end{proof}

\begin{corollary}[Cauchy 混合与同重心单核：全变差、相对熵与耗散的首项常数]\label{cor:xi-poisson-cauchy-mixture-point-leading-constants}
令 \(\tilde\nu\) 为 \(\RR\) 上概率测度，记其重心与方差为
\[
\bar\gamma:=\int_{\RR}\gamma\,\dd\tilde\nu(\gamma),
\qquad
\mathrm{Var}(\tilde\nu):=\int_{\RR}(\gamma-\bar\gamma)^2\,\dd\tilde\nu(\gamma).
\]
假设三阶中心绝对矩有限 \(\int_{\RR}|\gamma-\bar\gamma|^3\,\dd\tilde\nu(\gamma)<\infty\)，并定义
\[
h_t(x):=\int_{\RR}P_t(x-\gamma)\,\dd\tilde\nu(\gamma),
\qquad
g_t(x):=P_t(x-\bar\gamma).
\]
则当 \(t\to\infty\) 时成立锐主项
\begin{align}
\|h_t-g_t\|_{L^1(\RR)}&=\frac{3\sqrt3}{4\pi}\,\frac{\mathrm{Var}(\tilde\nu)}{t^2}+O(t^{-3}),\label{eq:xi-poisson-cauchy-mixture-tv-leading}\\
D_{\mathrm{KL}}(h_t\mid g_t)&=\frac{\mathrm{Var}(\tilde\nu)^2}{8t^4}+O(t^{-5}).\label{eq:xi-poisson-cauchy-mixture-kl-leading}
\end{align}
此外，令 \(\mathcal I_1\) 为 Poisson 流上的熵耗散率（命题 \ref{prop:xi-poisson-kl-dissipation-fractional-fisher}），则
\begin{equation}\label{eq:xi-poisson-cauchy-mixture-fisher-leading}
\mathcal I_1(h_t\mid g_t)=\frac{\mathrm{Var}(\tilde\nu)^2}{2t^5}+O(t^{-7}).
\end{equation}
从而得到单尺度反演公式
\begin{equation}\label{eq:xi-poisson-cauchy-mixture-variance-inversion}
\boxed{\;
\mathrm{Var}(\tilde\nu)
=\lim_{t\to\infty}\frac{4\pi}{3\sqrt3}\,t^2\|h_t-g_t\|_{L^1(\RR)}
=\lim_{t\to\infty}\sqrt{8}\,t^2\sqrt{D_{\mathrm{KL}}(h_t\mid g_t)}\; }.
\end{equation}
并且若存在 \(t_n\to\infty\) 使 \(t_n^4D_{\mathrm{KL}}(h_{t_n}\mid g_{t_n})\to 0\)，则 \(\mathrm{Var}(\tilde\nu)=0\)，从而 \(\tilde\nu=\delta_{\bar\gamma}\) 且 \(h_t\equiv g_t\)（\(\forall t>0\)）。
\end{corollary}

\begin{proof}
\eqref{eq:xi-poisson-cauchy-mixture-tv-leading} 由定理 \ref{thm:xi-cauchy-poisson-tv-sharp-constant} 给出；
\eqref{eq:xi-poisson-cauchy-mixture-kl-leading} 由定理 \ref{thm:xi-cauchy-poisson-kl-sharp-asymptotics}（取其首项）给出；
\par
其中常数 \(3\sqrt3/(4\pi)\) 与 \(1/8\) 的闭式评估可归结为一元积分
\[
\int_{\RR}\frac{|3y^2-1|}{(1+y^2)^3}\,\dd y=\frac{3\sqrt3}{4},
\qquad
\int_{\RR}\frac{(3y^2-1)^2}{(1+y^2)^5}\,\dd y=\frac{\pi}{4},
\]
参见命题 \ref{prop:cdim-poisson-kernel-derivative-l1-energy}（或等价地：定理 \ref{thm:xi-cauchy-poisson-tv-sharp-constant} 与定理 \ref{thm:xi-cauchy-poisson-kl-sharp-asymptotics} 的证明中的常数计算）。
\eqref{eq:xi-poisson-cauchy-mixture-fisher-leading} 由命题 \ref{prop:xi-poisson-kl-dissipation-fractional-fisher} 的耗散恒等式
\(
\frac{\dd}{\dd t}D_{\mathrm{KL}}(h_t\mid g_t)=-\mathcal I_1(h_t\mid g_t)
\)
与推论 \ref{cor:xi-cauchy-poisson-fisher-sharp-asymptotics} 的锐渐近给出。
两条反演公式由 \eqref{eq:xi-poisson-cauchy-mixture-tv-leading}--\eqref{eq:xi-poisson-cauchy-mixture-kl-leading} 直接同乘 \(t^2\) 或 \(t^4\) 并取极限得到。
最后一条刚性由 \eqref{eq:xi-poisson-cauchy-mixture-kl-leading} 或推论 \ref{cor:xi-cauchy-poisson-kl-rigidity} 立得：若 \(t_n^4D_{\mathrm{KL}}(h_{t_n}\mid g_{t_n})\to 0\)，则 \(\mathrm{Var}(\tilde\nu)^2/8=0\)。
而概率测度方差为零当且仅当其为 Dirac 测度，故 \(\tilde\nu=\delta_{\bar\gamma}\)，从而 \(h_t=P_t*\tilde\nu=g_t\)。证毕。
\end{proof}

\begin{remark}[可检索范围内的定位说明]\label{rem:xi-poisson-cauchy-mixture-search-note}
两条单一 Cauchy 分布之间的 \(D_{\mathrm{KL}}\) 闭式在公开资料中可检索；本段涉及的是 “Cauchy 混合 \(P_t*\tilde\nu\)” 与同重心单核 \(P_t(\cdot-\bar\gamma)\) 之间的信息损失首项常数与由此诱导的反演刚性。
以 “Poisson 卷积/粗粒化 + 相对熵 \(t^{-4}\) 主项常数” 等关键词组合在公开索引范围内作定位检索，未见与 \eqref{eq:xi-poisson-cauchy-mixture-kl-leading}--\eqref{eq:xi-poisson-cauchy-mixture-variance-inversion} 同型的直接陈述；
该类核验仅提供“可检索范围内未见同构”的接口，并不构成对未收录文献的穷尽性排除。
\end{remark}

\begin{conjecture}[端点化缺陷的二阶矩主导律：将一致趋零证书压缩为二阶矩一致匹配]\label{conj:xi-comoving-defect-second-moment-dominance}
设端点观测密度族 \(\{h^{(x_0)}\}_{x_0\in\RR}\) 及其粗粒化 \(h_t^{(x_0)}:=P_t*h^{(x_0)}\) 已按共动参数构造，并满足同均值规范化 \(\int x\,h^{(x_0)}(x)\dd x=0\)。
猜想存在一条参考密度 \(g\)（同均值且满足定理 \ref{thm:xi-poisson-cauchy-kl-fourth-order-universality} 的矩条件）与尺度对应 \(t=t(\varrho)\to\infty\)（当 \(\varrho\uparrow 1\)），使得在黄金分辨率链 \(\varrho=\varrho_m\) 上有
\[
\sup_{x_0\in\RR}
\Bigl|
D_{\zeta,x_0}(\varrho_m)-D_{\mathrm{KL}}\!\bigl(h^{(x_0)}_{t(\varrho_m)}\mid g_{t(\varrho_m)}\bigr)
\Bigr|
=o\!\bigl(t(\varrho_m)^{-4}\bigr),
\qquad
g_t:=P_t*g.
\]
若该猜想成立，则由定理 \ref{thm:xi-poisson-cauchy-kl-fourth-order-universality}--推论 \ref{cor:xi-poisson-cauchy-kl-moment-locking}，
\(\sup_{x_0}D_{\zeta,x_0}(\varrho_m)\to 0\) 的主导障碍可由二阶矩场
\[
\sup_{x_0\in\RR}\Bigl(\int_{\RR}x^2h^{(x_0)}(x)\dd x-\int_{\RR}x^2g(x)\dd x\Bigr)^2
\]
在 \(m\to\infty\) 时的归零控制给出；从而一致趋零证书在首项意义上被压缩为二阶矩的一致匹配问题。
\end{conjecture}

\begin{remark}[可检索范围内的定位说明]\label{rem:xi-comoving-transverse-window-literature-scope}
上述“横向可见窗—有限覆盖禁入”的证书形式是对共动缺陷上界的直接几何转译：其核心输入仅为 Cayley 模长闭式 \eqref{eq:xi-cayley-modulus-delta}、可见窗闭式 \eqref{eq:xi-single-defect-X} 与单零点 Jensen 下界（命题 \ref{prop:app-jensen-single-zero-lower-bound}）。
就这些关键词组合在公开可检索的文献与索引范围内作定位检索，未见与本文同一记号与证书口径下的直接同式陈述；该核验仅提供“可检索范围内未见同构”的否证接口，并不构成对未收录文献的穷尽性排除。
\end{remark}

\endinput
